\section{Mạng phản ứng và nhận biết chất}

\subsection{Không học từng phản ứng riêng lẻ}

Acid, base, oxide và muối liên hệ với nhau thành một mạng phản ứng. Khi đã nhìn được mạng này, bài chuỗi chuyển hóa và bài nhận biết trở nên ngắn hơn nhiều.

\begin{ghinho}{Mạng phản ứng cơ bản}
\[
\boxed{\text{oxide base}}
\longleftrightarrow
\boxed{\text{base}}
\longleftrightarrow
\boxed{\text{muối}}
\]
\[
\boxed{\text{oxide acid}}
\longleftrightarrow
\boxed{\text{acid}}
\longleftrightarrow
\boxed{\text{muối}}
\]
Các mũi tên không có nghĩa là mọi chất đều chuyển hóa trực tiếp theo mọi chiều. Mỗi bước phải kiểm tra xem có phản ứng phù hợp hay không.
\end{ghinho}

\subsection{Chuỗi chuyển hóa}

Khi gặp chuỗi $A\to B\to C$, hãy làm từng mũi tên riêng:
\begin{enumerate}
\item xác định loại chất ở đầu và cuối mũi tên;
\item chọn thuốc thử hoặc điều kiện phù hợp;
\item viết và cân bằng phương trình;
\item kiểm tra sản phẩm có đúng chất cần tạo hay không.
\end{enumerate}

\begin{vidu}
Hoàn thành chuỗi chuyển hóa:
\[
\ce{CaCO3 -> CaO -> Ca(OH)2 -> CaCO3}.
\]
\end{vidu}
\begin{loigiai}
\[
\ce{CaCO3 ->[t^\circ] CaO + CO2},
\]
\[
\ce{CaO + H2O -> Ca(OH)2},
\]
\[
\ce{Ca(OH)2 + CO2 -> CaCO3 v + H2O}.
\]

\textbf{Vậy,} chuỗi được thực hiện lần lượt bằng nhiệt phân, cho tác dụng với nước và dẫn \ce{CO2} vào dung dịch \ce{Ca(OH)2}.
\end{loigiai}

\subsection{Nhận biết bằng hiện tượng}

Một bài nhận biết tốt không bắt đầu bằng ``nhớ thuốc thử''. Hãy bắt đầu bằng câu hỏi: \textbf{các chất khác nhau ở tính chất nào có thể quan sát được?}

\begin{ghinho}{Hướng giải}
\begin{enumerate}
\item Liệt kê các chất cần phân biệt.
\item Tìm một thuốc thử làm xuất hiện các hiện tượng khác nhau.
\item Nếu chưa phân biệt hết, tiếp tục với các mẫu còn lại.
\item Ghi rõ hiện tượng và phương trình hóa học ở mỗi bước.
\end{enumerate}
\end{ghinho}

\begin{vidu}
Có ba dung dịch không nhãn: \ce{HCl}, \ce{NaOH}, \ce{NaCl}. Chỉ dùng quỳ tím, hãy nhận biết từng dung dịch.
\end{vidu}
\begin{loigiai}
Nhúng quỳ tím vào từng mẫu:
\begin{itemize}
\item mẫu làm quỳ tím hóa đỏ là \ce{HCl};
\item mẫu làm quỳ tím hóa xanh là \ce{NaOH};
\item mẫu không làm đổi màu quỳ tím là \ce{NaCl}.
\end{itemize}

\textbf{Vậy,} có thể nhận biết cả ba dung dịch chỉ bằng quỳ tím.
\end{loigiai}

\begin{tuluyen}
Hoàn thành chuỗi chuyển hóa sau và viết các phương trình hóa học:
\[
\ce{Cu -> CuO -> CuCl2 -> Cu(OH)2 -> CuO}.
\]
\dapso{\ce{2Cu + O2 -> 2CuO}; \ce{CuO + 2HCl -> CuCl2 + H2O}; \ce{CuCl2 + 2NaOH -> Cu(OH)2 v + 2NaCl}; \ce{Cu(OH)2 ->[t^\circ] CuO + H2O}.}
\begin{loigiaiBT}
Lần lượt thực hiện:
\[
\ce{2Cu + O2 ->[t^\circ] 2CuO},
\]
\[
\ce{CuO + 2HCl -> CuCl2 + H2O},
\]
\[
\ce{CuCl2 + 2NaOH -> Cu(OH)2 v + 2NaCl},
\]
\[
\ce{Cu(OH)2 ->[t^\circ] CuO + H2O}.
\]

\textbf{Vậy,} chuỗi chuyển hóa được hoàn thành bằng bốn phản ứng trên.
\end{loigiaiBT}
\end{tuluyen}
